% ============================================================
% CHAPTER 7: SRE METHODOLOGY & GAP ANALYSIS
% MCDA 5585 / 5586 Major Project Report
% Bhavik Kantilal Bhagat | A00494758 | Saint Mary's University
% ============================================================

\chapter{Methodologies \& Requirements}
\label{chap:methodologies}

During my internship I adopted Site Reliability Engineering as the primary
methodology, applying the four golden signals framework and SLI/SLO discipline
to the observability gap exposed by the CAIS-Pricing Silent Outage and the
Meridian Incident. This chapter covers the SRE principles I operationalised,
the gap analysis that established the starting posture, and the Dynatrace
parity evaluation that framed the research question.

% ────────────────────────────────────────────────────────────
\section{SRE Principles Applied}
\label{sec:sre_principles}
% ────────────────────────────────────────────────────────────

Site Reliability Engineering is an operational discipline that applies
software engineering thinking to the problem of keeping production
systems reliable at scale. Its conceptual vocabulary --- golden signals,
Service Level Indicators (SLIs), Service Level Objectives (SLOs), and
error budgets --- provided the theoretical backbone for every design
decision I made during this internship.

\subsection{The Four Golden Signals}

The four golden signals were introduced to the team by mentor Kristia
Marie Labos (Kri) during a dedicated knowledge-transfer session on
\textbf{May~25, 2026}. They define the minimum instrumentation required
to characterise the health of any service, regardless of the underlying
technology.

\begin{enumerate}
    \item \textbf{Latency} --- \emph{How long does a request take?} For
          the IA platform, latency manifests as Lambda function duration
          (measured at the P50, P95, and P99 percentiles), and as ECS
          service response time for API-facing containers. Sustained
          increases in P99 latency are an early warning of resource
          saturation or dependency degradation before errors appear.

    \item \textbf{Traffic} --- \emph{How much demand is on the system?}
          Traffic is measured as Lambda invocation count per unit time,
          SQS message ingestion rate for queue-backed automations, and
          MSK throughput (bytes in/out per second) for the Meridian
          Kafka pipeline. Traffic anomalies --- sudden spikes or
          unexpected drops --- often precede error and saturation
          signals.

    \item \textbf{Errors} --- \emph{What is the rate of failed
          requests?} The IA platform tracks errors as Lambda error and
          throttle counts, ECS task failure counts, and
          \texttt{NumberOfMessagesFailed} on SQS queues. An error rate
          that crosses a defined threshold indicates a code or
          infrastructure defect requiring immediate attention.

    \item \textbf{Saturation} --- \emph{How full is the service?}
          Saturation is the signal most relevant to the Apr~17,~2026
          Meridian incident. It is measured as Lambda concurrent
          execution count relative to the account concurrency limit,
          ECS task CPU and memory utilisation, MSK consumer lag (the
          number of messages awaiting processing on a Kafka topic), and
          SQS queue depth with oldest-message age. A service approaching
          saturation degrades before it fails --- making this the most
          actionable of the four signals.
\end{enumerate}

\subsection{Service Level Indicators and Objectives}

A \textbf{Service Level Indicator (SLI)} is a carefully defined
quantitative measure of some aspect of the level of service being
delivered. An SLI is not a raw metric --- it is the specific metric
chosen to represent one of the four golden signals for a particular
service. A \textbf{Service Level Objective (SLO)} is a target value
or range of values for an SLI, measured over a rolling time window.
Together, SLIs and SLOs transform vague aspirations such as
``the system should be reliable'' into falsifiable engineering
commitments.

For the IA platform, I adopted the following SLOs as the design
targets for the observability build:

\begin{itemize}
    \item \textbf{Lambda error rate SLO}: Fewer than 1\% of Lambda
          invocations must result in an error over any 24-hour rolling
          window. This translates directly to the CloudWatch alarm
          threshold applied to the \texttt{Errors} metric with
          \texttt{ErrorRate = Errors / Invocations}.

    \item \textbf{MSK consumer lag SLO}: The Kafka consumer lag for
          any Meridian consumer group must remain below
          \textbf{10,000~messages} at any point in time. This threshold
          was derived from the Apr~17 post-mortem: the system reached
          9~million messages in backlog before the business team
          escalated, and a threshold of 10,000 would have triggered an
          alarm within minutes of the first deviation.

    \item \textbf{ECS CPU utilisation SLO}: No ECS service should
          sustain CPU utilisation above \textbf{80\%} over a 5-minute
          evaluation period. Breaching this threshold activates the
          auto-scaling policy and triggers an SNS notification to the
          SRE team.

    \item \textbf{SQS queue depth SLO}: No production SQS queue should
          accumulate more than a configurable maximum of unprocessed
          messages before an alarm fires. The threshold is set per-queue
          based on the expected processing throughput for that
          automation.
\end{itemize}

\subsection{Dynatrace Parity Analysis}
\label{sec:dynatrace_parity}

A key framing question for my internship --- introduced in
Chapter~\ref{chap:executive_summary} and developed here in methodological
detail --- is how closely an AWS-native observability stack can replicate
the capabilities of Dynatrace, and at what incremental cost. The PoC
investigation evaluated 25 Dynatrace features against five AWS tool
combinations to quantify achievable parity.

Table~\ref{tab:dynatrace_parity} summarises the parity analysis for the
ten most operationally relevant Dynatrace capabilities to the IA platform.

\begin{table}[ht]
\centering
\caption{Dynatrace Capability Parity --- AWS-Native Stack vs Dynatrace (selected features)}
\label{tab:dynatrace_parity}
\begin{tabular}{p{4.5cm}p{1.4cm}p{1.4cm}p{1.4cm}p{1.8cm}p{1.4cm}}
\toprule
\textbf{Dynatrace Feature} & \textbf{CW} & \textbf{+Grafana} & \textbf{+X-Ray} & \textbf{+App Signals} & \textbf{+Tempo} \\
\midrule
Alerting                  & \checkmark & \checkmark & --- & --- & --- \\
Infrastructure Metrics    & \checkmark & \checkmark & --- & --- & --- \\
Log Analytics             & \checkmark & \checkmark & --- & --- & --- \\
Container Monitoring      & \checkmark & \checkmark & --- & \checkmark & --- \\
Error Rate Tracking       & \checkmark & \checkmark & \checkmark & \checkmark & --- \\
Latency Percentiles       & --- & \checkmark & \checkmark & \checkmark & \checkmark \\
Distributed Trace Waterfall & --- & --- & \checkmark & --- & \checkmark \\
Service Map / Topology    & --- & --- & \checkmark & \checkmark & --- \\
SLO Management            & --- & --- & --- & \checkmark & --- \\
Dependency Mapping        & --- & --- & \checkmark & \checkmark & --- \\
\midrule
\textbf{Parity (\% of 25 features)} & 30\% & 33\% & 60\% & 80\% & 83\% \\
\textbf{Monthly Cost (est.)} & \$49 & \$68 & \$77 & \$232 & \$269 \\
\bottomrule
\end{tabular}
\end{table}

\noindent
Three capabilities remain \textbf{permanently outside the AWS-native reach}:
Davis~AI causal analysis (proprietary ML, no AWS equivalent); automatic
baselining (Dynatrace learns per-service baselines over weeks; AWS alarms
require manual thresholds); and session replay (full DOM capture, not
available in any OSS tool). These account for approximately 9--10\% of
the functional gap that no combination of AWS tools can close.

\paragraph{Adopted stack.}
The PoC I delivered in this internship corresponds to the
\textbf{CloudWatch + Grafana + X-Ray} column: 60\% parity at approximately
\$77/month. Application Signals (the path to 80\% parity) requires deploying
an ADOT sidecar on each ECS service and is scoped as post-internship work.
The gap-to-tooling mapping is used throughout Chapter~\ref{chap:implementation}
and Chapter~\ref{chap:achievements} to contextualise what the PoC does and does
not cover.

\subsection{Error Budgets}

An \textbf{error budget} is the complement of the SLO: it is the
maximum permissible unreliability for a service over a given window.
If the SLO is that Lambda error rate must be below 1\% over 24~hours,
then the error budget is the 1\% of invocations that are permitted to
fail before the SLO is breached. When the error budget is consumed ---
either by a single outage or by sustained degradation --- the SRE
engineering principle dictates that all non-reliability work (feature
development, infrastructure experiments) should be paused in favour of
reliability remediation.

In the context of my IA internship, I applied error budgets
\emph{conceptually} rather than through automated budget-burn tracking.
The Grafana dashboard exposes the current SLI values for each monitored
service, and the CloudWatch alarms fire when the system approaches the
SLO threshold --- effectively providing a manual signal that the error
budget is being drawn down. Full automated budget-burn rate alerting
(which would require custom CloudWatch Metric Math expressions computing
burn rate against a rolling budget) was identified as a post-internship
enhancement.

% ────────────────────────────────────────────────────────────
\section{Observability Gap Analysis}
\label{sec:observability_gap}
% ────────────────────────────────────────────────────────────

A prerequisite for building any monitoring platform is understanding the
current state of observability. Before I wrote a single line of code
for IISS-487, I performed a structured gap analysis to characterise
what the IA platform could and could not see about itself.

\subsection{Pre-Internship Monitoring Posture}

At the commencement of my internship in Mar~2026, the IA-SRE team's
monitoring posture exhibited the following deficiencies:

\begin{itemize}
    \item \textbf{No CloudWatch alarms on any Meridian service.} The
          Meridian Kafka-backed event pipeline --- the most
          business-critical component of the IA platform --- had zero
          proactive alerting. No alarms existed on MSK consumer lag,
          ECS task failure counts, Lambda error rates, or SQS queue
          depth for any Meridian-owned resource.

    \item \textbf{No centralized dashboard.} Monitoring was entirely
          fragmented. Operators who wished to check the health of a
          service were required to navigate to that service's individual
          CloudWatch console page. There was no eagle-eye view
          from which to assess the overall health of the platform.

    \item \textbf{Incomplete Container Insights coverage.} Of the six
          Meridian ECS clusters, only one
          (\texttt{meridian-compute-cluster}, enabled since Jul~2024)
          had Amazon CloudWatch Container Insights activated. The
          remaining five clusters were invisible to task-level CPU and
          memory telemetry.

    \item \textbf{No MSK consumer lag monitoring.} Despite Kafka
          consumer lag being the primary leading indicator of Meridian
          processing health --- as demonstrated conclusively by the
          Meridian Incident --- no metric or alarm existed for this
          signal. The \texttt{EstimatedMaxTimeLag} and
          \texttt{SumOffsetLag} metrics available natively from Amazon
          MSK were not surfaced anywhere.

    \item \textbf{No SLO tracking or error budget calculation.}
          The team had no defined SLIs, no published SLO targets, and
          therefore no basis for calculating error budgets or
          determining when the platform was in an unacceptable
          reliability state.

    \item \textbf{No automated resource inventory.} The inventory of
          Lambda functions, ECS clusters, and EC2 instances was
          maintained manually via Confluence pages and spreadsheets.
          This made it impossible to know, at any given moment, exactly
          which resources existed across all six platforms.

    \item \textbf{RPA bot monitoring limited to orchestrator UIs.}
          Blue Prism and UiPath bots were monitored exclusively through
          their respective orchestrator platforms (Blue Prism Orchestrator and
          UiPath Cloud). No CloudWatch integration existed, and there
          was no mechanism for the SRE team to receive proactive
          notification of bot failures or queue backlogs outside the
          orchestrator UI.

    \item \textbf{Reactive operational posture.} The cumulative effect
          of the above gaps was that the team's primary incident
          detection mechanism was business-team escalation. When a fund
          processing automation failed silently, the first signal
          received by the SRE team was a ticket or message from
          operations staff in Dublin, London, or Halifax who had noticed
          anomalous fund processing outputs. This reactive posture
          introduced structural latency between failure onset and
          engineering response.
\end{itemize}

\subsection{Gap Analysis Summary}

Table~\ref{tab:gap_analysis} summarises the pre-internship gaps and the
post-internship resolutions for each monitoring domain.

\begin{table}[ht]
\centering
\caption{Observability Gap Analysis --- Pre- vs.\ Post-Internship}
\label{tab:gap_analysis}
\begin{tabular}{p{3.5cm}p{4.5cm}p{4.5cm}}
\toprule
\textbf{Domain} & \textbf{Pre-Internship State} & \textbf{Post-Internship State} \\
\midrule
CloudWatch Alarms & 0 alarms on any IA resource & 38 alarms across 7 AWS service categories \\
\midrule
Centralized Dashboard & None & 14~rows, 110+~panels (Grafana AMG) \\
\midrule
Container Insights & 1 of 6 Meridian clusters & 5 of 6 Meridian clusters \\
\midrule
MSK Consumer Lag & No monitoring & CloudWatch alarm on \texttt{EstimatedMaxTimeLag} ($<$30~sec detection) \\
\midrule
SLO Tracking & None & SLI/SLO targets defined; dashboard exposes current values \\
\midrule
Resource Inventory & Manual (Confluence) & Automated REST API (Lambda, ECS, EC2 across all platforms) \\
\midrule
Incident Detection & Business team escalation & Proactive alarm with SNS routing \\
\bottomrule
\end{tabular}
\end{table}

% ────────────────────────────────────────────────────────────
\section{Limitations and Mitigations}
\label{sec:limitations}
% ────────────────────────────────────────────────────────────

Several CloudWatch and AMG constraints shaped the final architecture and
required deliberate design decisions to work around.

\textbf{CloudWatch \texttt{SEARCH()} visibility gap.} The \texttt{SEARCH()}
function only returns metrics for AWS services with recent activity. Lambda
functions not invoked within the lookback window disappear from the
dashboard entirely --- the same root cause that made the CAIS pipeline
invisible in IISS-430. \textit{Mitigation:} I replaced all
\texttt{SEARCH()} expressions with CloudWatch Metric Insights SQL
(\texttt{SELECT ... GROUP BY FunctionName}), which enumerates all
deployed AWS services in the namespace regardless of recent activity.

\textbf{Metric Insights SQL cannot compute percentiles.} The
\texttt{GROUP BY} mode of Metric Insights only supports \texttt{AVG},
\texttt{SUM}, \texttt{MIN}, and \texttt{MAX}. P50/P99 latency statistics
are unavailable in this mode. \textit{Mitigation:} Latency percentile
panels use CloudWatch standard metrics mode with a
\texttt{FunctionName:~["*"]} wildcard dimension, which restores true
percentiles but reintroduces the visibility gap for idle AWS services ---
an accepted tradeoff documented in the dashboard design.

\textbf{CloudWatch subscription filters blocked by IAM boundary.}
A proposed central log-aggregation design would require 106 metric
subscription filters across the six platforms. Analysis confirmed this
was technically feasible but blocked by the
\texttt{logs:PutSubscriptionFilter} action not being granted within the
team's scoped permission set. \textit{Mitigation:} I adopted CloudWatch
Logs Insights widgets embedded in Grafana panels for log-based error
surfacing, avoiding subscription filters entirely.

\textbf{Secrets Manager 64\,KB payload limit.} The Lambda-backed
CloudFormation provisioner stores alarm templates as a Secrets Manager
payload. As the monitoring scope expanded to cover all AWS service
categories, the generated YAML exceeded the hard 64\,KB limit.
\textit{Mitigation:} I migrated all templates to an S3 bucket with
object versioning; the provisioner fetches templates via
\texttt{s3:GetObject} at deploy time.

\textbf{Shared \texttt{BudgetCode} tag across platforms.} The
\texttt{meridian} and \texttt{citcoworks} platforms share the internal
tag value \texttt{"Shared Cloud Artifacts"}, making a single-stage tag
lookup ambiguous when discovering AWS services in those platforms.
\textit{Mitigation:} I designated \texttt{AppName} as the mandatory
secondary discriminator for all AWS service discovery in those platforms,
documented as a platform-wide governance constraint in the inventory API
\texttt{README}.
